% 原文 pp.122–127（印刷 pp.1240–1245）

% ===== 原文 p.122 =====
% 承接 p.121 末句（"若在深处观测 $\phi$……"）
在 $C$ 内部的深处，$\phi$ 只作为 $x$ 的函数。于是需要极小化的是
\begin{equation}
  \tag{8.57}
  \label{eq:8-57}
  \int dx\left(\left(\frac{d\phi}{dx}\right)^{2}-2e^{-I/e^{2}}\left(\cos 2\phi-1\right)\right),
\end{equation}
边界条件为：当 $x\to-\infty$ 时 $\phi\to0$，当 $x\to+\infty$ 时 $\phi\to\pi$。
这个带边界条件的变分问题有一个解，且该解在 $x$ 的平移下是唯一的；由于存在质量间隙，$\phi$ 以指数快的速度趋近其渐近值。
记这个一维解的作用量为 $S$。

% [Tutor Note: (8.57) 只是示意性的写法，不能与 (8.56) 逐字对照。把 (8.56) 按 $\phi=\phi(x)$ 约化（另两个方向取单位面积），逐项系数应为 $\frac{1}{4\pi e^{2}}\int dx\big[(\phi')^{2}-8\pi e^{2}e^{-I/e^{2}}(\cos2\phi-1)\big]$；而 (8.57) 相当于把整体因子 $4\pi e^{2}$ 吸收掉、势的系数却仍写作 $2e^{-I/e^{2}}$，两者相差一个因子 $4\pi e^{2}$。也就是说，(8.57) 隐含地对 $x$ 作了一次标度变换并对 $S$ 重新定了标。kink 解的形状、其指数衰减率以及面积律的结论都不受影响，但若照 (8.57) 直接计算 $S$，就会得到关于 $e$ 的错误幂次——弦张力 $\sigma\sim\frac{1}{e}e^{-I/2e^{2}}$ 中的 $1/e$ 只能来自被作者略去的那个整体因子 $\frac{1}{4\pi e^{2}}$。另外，「指数快地趋近渐近值」的依据是势在 $\phi=0,\pi$ 附近的曲率 $m^{2}\propto e^{-I/e^{2}}$，即质量间隙；正因 $m$ 指数地小，畴壁（domain wall）才是「薄」的。]

现在考虑一般的情形：若 $D$ 是以 $C$ 为边界的极小面积曲面，$A(D)$ 是它的面积，那么作为相应平坦丛的截面、且使作用量取极小的那个 $\phi$ 具有如下性质：除 $D$ 附近之外，$\phi$ 都非常接近 $0$ 或 $\pi$；穿过 $D$ 时它跳跃 $\pi$；而在 $D$ 附近，它的剖面看起来正像上一段所讨论的理想化一维问题的解。
它的作用量非常接近 $A(D)S$，因此 Wilson 线的期望值近似为 $e^{-A(D)S}$，这就显示出面积律与禁闭。

\subsection{四维中的对偶与 $SL(2,\mathbb{Z})$}
% 原文 §8.5：Duality in four dimensions and $SL(2,\mathbb{Z})$

接下来我们讨论四维中的阿贝尔对偶（abelian duality），并把它应用于四维自由 $U(1)$ 理论的 $SL(2,\mathbb{Z})$ 对称性。
非自由理论的 $SL(2,\mathbb{Z})$ 对称性，我们留到后面某一讲再讨论，因为迄今所知的所有这类例子都涉及超对称。

我们在一个四维流形 $M$ 上取一个 $U(1)$ 丛 $\mathcal{L}$，以及 $\mathcal{L}$ 上的一个联络 $A$，其曲率为 $F=F_A$。
包含拓扑项的规范理论拉氏量（欧氏符号约定，Euclidean signature）为
\begin{equation}
  \tag{8.58}
  \label{eq:8-58}
  \begin{aligned}
  \mathcal{L}(A)&=\frac{1}{4e^{2}}\int d^{4}x\,\sqrt{g}\,F_{mn}F^{mn}+\frac{i\theta}{16\pi^{2}}\int d^{4}x\,\sqrt{g}\,\epsilon_{mnpq}F^{mn}F^{pq}\\
  &=\frac{1}{2e^{2}}\int F_{A}\wedge *F_{A}+\frac{i\theta}{4\pi^{2}}\int F_{A}\wedge F_{A}.
  \end{aligned}
\end{equation}

我们在动能项上采用了标准的归一化，并把拓扑项归一化，使得把 $\theta$ 换为 $\theta+2\pi$ 不改变物理。
（拓扑项的这一性质来自这样一个事实：$c_{1}(\mathcal{L})^{2}=\int(F_{A}/2\pi)\wedge(F_{A}/2\pi)$ 恒为整数。注意，在自旋流形上 $c_{1}(\mathcal{L})^{2}$ 恒为\emph{偶}整数，于是我们又多出一个等价性：把 $\theta$ 换为 $\theta+\pi$ 也不改变物理。）

令 $\tau=\frac{\theta}{\pi}+\frac{2\pi i}{e^{2}}\in\mathfrak{h}$（$\mathfrak{h}$ 为复上半平面）。
正如我们刚才注意到的，$\tau\mapsto\tau+2$ 是这个理论的一个对称性；而在自旋流形上工作时，$\tau\mapsto\tau+1$ 也是对称性。
要把它扩充为一个 $SL(2,\mathbb{Z})$ 作用（在自旋流形的情形），我们还需要一个把 $\tau$ 映为 $-1/\tau$ 的对称性；这将由对偶变换 $F_{A}\leftrightarrow *F_{C}$ 给出（其中 $C$ 是一个新的「对偶」联络）。

这一对偶变换的计算与 §8.2、§8.3 中的计算类似。
我们先定义 $F_{\pm}=\frac12(F_{A}\pm *F_{A})$，并把拉氏量 (8.58) 改写为
\begin{equation}
  \tag{8.59}
  \label{eq:8-59}
  \begin{aligned}
  \mathcal{L}(A)&=\frac{i\bar\tau}{4\pi}\int F_{+}\wedge F_{+}+\frac{i\tau}{4\pi}\int F_{-}\wedge F_{-}\\
  &=\frac{i\bar\tau}{4\pi}\int\|F_{+}\|^{2}-\frac{i\tau}{4\pi}\int\|F_{-}\|^{2}.
  \end{aligned}
\end{equation}

% [Tutor Note: 由 (8.58) 到 (8.59) 这两步值得逐项验算，因为 $\tau$ 定义中 $\frac{2\pi}{e^{2}}$（而不是 $\frac{\pi}{e^{2}}$）最容易记错。由 $F_{\pm}=\frac12(F_{A}\pm *F_{A})$ 及自对偶/反自对偶性 $*F_{\pm}=\pm F_{\pm}$ 得 $F_{A}=F_{+}+F_{-}$、$*F_{A}=F_{+}-F_{-}$，于是 $F_{A}\wedge *F_{A}=\|F_{+}\|^{2}+\|F_{-}\|^{2}$、$F_{A}\wedge F_{A}=\|F_{+}\|^{2}-\|F_{-}\|^{2}$（交叉项 $F_{+}\wedge F_{-}$ 因次数不相容而消失）。代入 (8.58) 得 $\|F_{+}\|^{2}$ 的系数 $\frac{1}{2e^{2}}+\frac{i\theta}{4\pi^{2}}$、$\|F_{-}\|^{2}$ 的系数 $\frac{1}{2e^{2}}-\frac{i\theta}{4\pi^{2}}$；而 $\frac{i\bar\tau}{4\pi}=\frac{i}{4\pi}\big(\frac{\theta}{\pi}-\frac{2\pi i}{e^{2}}\big)=\frac{1}{2e^{2}}+\frac{i\theta}{4\pi^{2}}$、$-\frac{i\tau}{4\pi}=\frac{1}{2e^{2}}-\frac{i\theta}{4\pi^{2}}$，恰好吻合。第二行那个负号来自 $F_{-}\wedge F_{-}=-\|F_{-}\|^{2}$——反自对偶二形式的外积为负，这是欧氏四维中最常被漏掉的一个符号；若把 $F_{-}\wedge F_{-}$ 也当成 $+\|F_{-}\|^{2}$，就会得出 $\tau$ 与 $\bar\tau$ 位置互换的错误结论，进而把下面模形式的权全部算错。]

% ===== 原文 p.123 =====
记 $\mathcal{G}$ 为与 $A$ 相联系的规范群，则这个理论的配分函数可以写成
\begin{equation}
  \tag{8.60}
  \label{eq:8-60}
  Z(\tau)=\frac{1}{\mathrm{vol}(\mathcal{G})}\sum_{\mathcal{L}}\int\mathcal{D}A\,\exp\left(-\frac{i\bar\tau}{4\pi}\int\|F_{+}\|^{2}+\frac{i\tau}{4\pi}\int\|F_{-}\|^{2}\right).
\end{equation}

我们早先的对偶例子都从这样一个标量场 $\phi$ 的理论出发：$\phi$ 只通过其导数 $d\phi$ 进入拉氏量，因而理论具有 $\phi\mapsto\phi+c$（$c$ 为常数）的对称性；对偶变换的第一步是把这一对称性规范化（gauge），同时引入一个适当的 Lagrange 乘子场。

而现在的理论本身已经是一个规范理论：它是联络 $A$ 的理论，$A$ 只通过其曲率 $F_{A}$ 进入拉氏量，于是理论具有 $A\mapsto A+B$ 的对称性（$B$ 为平坦联络）。
我们要做的，就是「把这一对称性规范化」的类比：允许 $B$ 是任意丛上的任意联络，并引入一种「奇异规范场」（exotic gauge field）$G$——它是一个 2-形式场——把对称性扩充为
\begin{equation}
  \tag{8.61}
  \label{eq:8-61}
  \begin{aligned}
  A&\to A+B\\
  G&\to G+F_{B}.
  \end{aligned}
\end{equation}
于是 $\mathcal{F}:=F_{A}-G$ 扮演「规范不变」量的角色，类似于标量场的协变导数。
需要强调的是：两个 $G$ 场若相差 $G\to G+F_{B}$（$F_{B}$ 为任意线丛上任意联络的曲率），就认为它们是规范等价的。
在我们的分析中，为简单起见，我们假定 $H^{2}(M)$ 没有挠（torsion）。

我们需要拉氏量的一个「规范不变」的扩充。我们或许可以尝试
\begin{equation}
  \tag{8.62}
  \label{eq:8-62}
  \mathcal{L}(A,G)=\frac{i\bar\tau}{4\pi}\int\|\mathcal{F}_{+}\|^{2}-\frac{i\tau}{4\pi}\int\|\mathcal{F}_{-}\|^{2},
\end{equation}
但这太简单了（因为，举例来说，其经典方程给出 $\mathcal{F}=0$）。
为改进这一点，我们在一个新的线丛 $\mathcal{N}$ 上引入一个新的联络 $C$，其曲率为 $F_{C}$，并考虑拉氏量
\begin{equation}
  \tag{8.63}
  \label{eq:8-63}
  \mathcal{L}(A,G,C)=\frac{i\bar\tau}{4\pi}\int\|\mathcal{F}_{+}\|^{2}-\frac{i\tau}{4\pi}\int\|\mathcal{F}_{-}\|^{2}-\frac{i}{2\pi}\int F_{C}\wedge G.
\end{equation}

这个新理论的配分函数可以表示为一个路径积分，其中包含与丛 $\mathcal{L}$ 和 $\mathcal{N}$ 的一切可能选取相对应的各个扇区：
\begin{equation}
  \tag{8.64}
  \label{eq:8-64}
  \begin{aligned}
  &\frac{1}{\mathrm{vol}(\widetilde{\mathcal{G}})}\,\frac{1}{\mathrm{vol}(\mathcal{G})}\,\frac{1}{\mathrm{vol}(\mathcal{G}_{C})}\sum_{\mathcal{L},\mathcal{N}}\int\mathcal{D}A\,\mathcal{D}G\,\mathcal{D}C\\
  &\qquad\times\exp\left(-\frac{i\bar\tau}{4\pi}\int\|\mathcal{F}_{+}\|^{2}+\frac{i\tau}{4\pi}\int\|\mathcal{F}_{-}\|^{2}+\frac{i}{2\pi}\int F_{C}\wedge G\right),
  \end{aligned}
\end{equation}
其中 $\mathcal{G}$ 与 $\mathcal{G}_{C}$ 分别表示与 $A$ 和 $C$ 相联系的规范群，而 $\widetilde{\mathcal{G}}$ 表示「奇异」规范群。

为看出这个新理论与原理论等价，我们先在 (8.64) 中作 $C$ 积分：把 $C$ 写成 $C=C_{0}+C'$，其中 $C_{0}$ 是线丛 $\mathcal{N}$ 上取定的一个联络。于是对 $C'$ 的积分是
\begin{equation}
  \tag{8.65}
  \label{eq:8-65}
  \frac{1}{\mathrm{vol}\,\mathcal{G}_{C}}\int\mathcal{D}C'\,e^{\frac{i}{2\pi}\int C'\wedge dG}=\delta(dG).
\end{equation}

% [Tutor Note: (8.65) 之所以只剩 $\delta(dG)$，关键在于 $\int F_{C'}\wedge G=\int C'\wedge dG$（分部积分，$dG$ 为闭形式时边界项消失），因此这一积分与 $\mathcal{N}$ 的选取、与 $C_{0}$ 都无关，只剩一个纯粹约束 $dG=0$；也正因如此，后面 (8.66) 才能把它与对 $\mathcal{N}$ 的求和分开处理。另一个易错点是：$G$ 是 2-形式，$dG$ 是 3-形式，而 $M$ 是四维的，故 $\delta(dG)$ 是 3-形式空间上的泛函 delta 函数，其「体积」归一化正是 (8.64) 里 $\frac{1}{\mathrm{vol}(\widetilde{\mathcal{G}})}$ 这个因子的来源。若把 $G$ 误当作 1-形式（像通常的 $U(1)$ 对偶那样），就会把 (8.63) 的 $F_{C}\wedge G$ 误写成 $F_{C}\wedge *G$，整个推导随之垮掉。]

% ===== 原文 p.124 =====
于是，当我们对 $\mathcal{N}$ 求和时，得到
\begin{equation}
  \tag{8.66}
  \label{eq:8-66}
  \frac{1}{\mathrm{vol}\,\mathcal{G}_{C}}\sum_{\mathcal{N}}\int\mathcal{D}C\,e^{-\frac{i}{2\pi}\int F_{C}\wedge G}=\sum_{x\in H^{2}(M)}e^{i(x,G)}\,\delta(dG)=\delta\left(\left[\frac{G}{2\pi}\right]\in\mathbb{Z}\right)\delta(dG).
\end{equation}

$dG=0$ 以及 $[G/2\pi]$ 是整上同调类这两个条件，恰恰意味着 $G$ 具有 $F_{B}$ 的形式（$F_{B}$ 是某个线丛上某个联络的曲率），因而意味着 $G$ 可以被规范到零。
做完这一步之后，就可推出配分函数 (8.64) 与 $Z(\tau)$ 一致，于是我们回到了原来的理论。

换一种做法：我们也可以利用「奇异」规范不变性把 $A$ 规范到 $0$，从而计算配分函数 (8.64)。这样剩下的路径积分是
\begin{equation}
  \tag{8.67}
  \label{eq:8-67}
  \frac{1}{\mathrm{vol}\,\mathcal{G}_{C}}\sum_{\mathcal{N}}\int\mathcal{D}G\int\mathcal{D}C\,\exp\left(-\frac{i\bar\tau}{4\pi}\int\|G_{+}\|^{2}+\frac{i\tau}{4\pi}\int\|G_{-}\|^{2}+\frac{i}{2\pi}\int F_{C}\wedge G\right).
\end{equation}

为计算 $G$ 积分，我们配方，同时记住
\begin{equation}
  \tag{8.68}
  \label{eq:8-68}
  \int F_{C}\wedge G=\int\left(F_{C+}\wedge *G_{+}-F_{C-}\wedge *G_{-}\right).
\end{equation}
事实上，若定义 $G'=G-\frac{1}{\bar\tau}F_{C+}+\frac{1}{\tau}F_{C-}$，就可以把 eq. (8.67) 中的指数写成
\begin{equation}
  \tag{8.69}
  \label{eq:8-69}
  -\frac{i\bar\tau}{4\pi}\int\|G'_{+}\|^{2}+\frac{i\tau}{4\pi}\int\|G'_{-}\|^{2}+\frac{i}{4\pi\bar\tau}\int\|F_{C+}\|^{2}-\frac{i}{4\pi\tau}\int\|F_{C-}\|^{2}.
\end{equation}

当我们完成 $G'$ 积分时，前两项给出一个高斯积分，它只对整体归一化有贡献；把 $G'$ 积掉之后，剩下的路径积分是
\begin{equation}
  \tag{8.70}
  \label{eq:8-70}
  \frac{1}{\mathrm{vol}\,\mathcal{G}_{C}}\sum_{\mathcal{N}}\int\mathcal{D}C\,\exp\left(\frac{i}{4\pi\bar\tau}\int\|F_{C+}\|^{2}-\frac{i}{4\pi\tau}\int\|F_{C-}\|^{2}\right).
\end{equation}

这与原来的路径积分相同，只是 $\tau$ 被换成了 $-1/\tau$——这正是我们想要证明的。

% [Tutor Note: (8.69) 的配方建议自己动手复算，因为 $G'$ 的平移量中 $\frac{1}{\bar\tau}$ 与 $\frac{1}{\tau}$ 的搭配不对称，极易写反。以 $G_{+}$ 部分为例，指数为 $-\frac{i\bar\tau}{4\pi}\|G_{+}\|^{2}+\frac{i}{2\pi}\langle F_{C+},G_{+}\rangle$，形如 $a\|G\|^{2}+b\langle F,G\rangle=a\big\|G+\frac{b}{2a}F\big\|^{2}-\frac{b^{2}}{4a}\|F\|^{2}$，此处 $a=-\frac{i\bar\tau}{4\pi}$、$b=\frac{i}{2\pi}$，故 $\frac{b}{2a}=-\frac{1}{\bar\tau}$、$-\frac{b^{2}}{4a}=+\frac{i}{4\pi\bar\tau}$；$G_{-}$ 部分同法给出平移 $-\frac{1}{\tau}F_{C-}$ 与残余项 $-\frac{i}{4\pi\tau}$。四个系数与 (8.69) 逐字相符，可见原文此处无误。真正的物理在残余项：$\frac{i}{4\pi\bar\tau}$ 与 $-\frac{i}{4\pi\tau}$ 恰好是把原作用量中 $\|F_{C+}\|^{2}$、$\|F_{C-}\|^{2}$ 的系数 $\frac{i\bar\tau}{4\pi}$、$-\frac{i\tau}{4\pi}$ 里的 $\tau$ 换成 $-1/\tau$（注意 $\overline{-1/\tau}=-1/\bar\tau$）后所得的结果——高斯积分把 $\tau$ 翻转成 $-1/\tau$，这就是 $S$ 对偶的全部机制，而不是某个额外假设。]

正如我们在两维情形所做的那样，也可以分析路径积分归一化的 $\tau$ 依赖性，从而得到更多有趣的结果。
若计算对 $G'$ 的高斯积分，就能窥见用这种方式所得结果之风味的端倪，该积分给出
\begin{equation}
  \tag{8.71}
  \label{eq:8-71}
  \left(\frac{2\pi}{\sqrt{i\bar\tau}}\right)^{n_{2+}}\left(\frac{2\pi}{\sqrt{-i\tau}}\right)^{n_{2-}},
\end{equation}
其中 $n_{2\pm}$ 表示自对偶与反自对偶 2-形式的个数。
当然，这两个数都是无穷大，所以必定存在与其他归一化因子之间的某种相消。
把这一点详细算出来之后，\footnote{见参考文献中所引的我的论文。}所得的结果是
\begin{equation}
  \tag{8.72}
  \label{eq:8-72}
  Z(\tau)=\tau^{-(\chi+\sigma)/4}\,\bar\tau^{-(\chi-\sigma)/4}\,Z(-1/\tau),
\end{equation}
其中 $\chi$ 与 $\sigma$ 分别是 $M$ 的欧拉数与号差。
于是，配分函数 $Z(\tau)$ 实际上是 $SL(2,\mathbb{Z})$（当流形非自旋时，是某个子群）的一个权为 $\left(\frac{\chi+\sigma}{4},\frac{\chi-\sigma}{4}\right)$ 的模形式（modular form）。

% [Tutor Note: (8.71) 的指数 $n_{2\pm}$ 是自对偶/反自对偶调和二形式的个数：对紧流形它们有限，但一般依赖于度规；而 (8.72) 把幂次换成了 $\frac{\chi+\sigma}{4}$ 与 $\frac{\chi-\sigma}{4}$，其间做了一次正规化。按热核/$\zeta$ 函数正规化，行列式比值的有限部分给出 $n_{2+}\to\frac{\chi+\sigma}{2}$、$n_{2-}\to\frac{\chi-\sigma}{2}$（注意 $\frac{\chi+\sigma}{2}=1-b_{1}+b_{2}^{+}$，正是 de Rham 复形在偶形式上的指标），这就是 Atiyah--Singer 指标定理中 Gauss--Bonnet 与号差两个示性类组合的体现。作者在正文只说「必定存在与其他归一化因子之间的相消」，细节见本讲末参考文献中 E. Witten, \emph{On $S$-duality in abelian gauge theory}, Selecta Math. (N.S.) 1 (1995)。提醒：从 (8.71) 直接读出的权是 $(\frac{n_{2+}}{2},\frac{n_{2-}}{2})$，与 (8.72) 的 $(\frac{\chi+\sigma}{4},\frac{\chi-\sigma}{4})$ 并不相同，两者之差正是上述正规化，切不可混用。]

% ===== 原文 p.125 =====
我们也可以像在低维情形中那样，让某些算符插入穿过对偶变换。
原理论中 $F_{\pm}$ 的插入，在扩充理论中可以由插入 $\mathcal{F}_{\pm}=F_{\pm}-G$ 来实现，而在把变量换成 $G'$ 之后，这可以分别写成
\begin{equation}
  \tag{8.73}
  \label{eq:8-73}
  \mathcal{F}_{+}=F_{+}-G'_{+}-\frac{1}{\bar\tau}F_{C+},\qquad\text{或}\qquad \mathcal{F}_{-}=F_{-}-G'_{-}+\frac{1}{\tau}F_{C-},
\end{equation}
于是，当我们把 $A$ 规范到零并积掉 $G'$ 之后，就只剩下正比于 $F_{C\pm}$ 的算符插入，即：
\begin{equation}
  \tag{8.74}
  \label{eq:8-74}
  F_{+}\mapsto\left(-1/\bar\tau\right)F_{C+},\qquad\text{以及}\qquad F_{-}\mapsto\left(1/\tau\right)F_{C-}.
\end{equation}

请注意：由于这些映射带有 $\tau$ 依赖性，一个含有 $F_{+}$ 与 $F_{-}$ 插入的关联函数，其模权将与配分函数的模权不同。

\subsection{Hamilton 形式}
% 原文 §8.6：The Hamiltonian formalism

回到规范群为 $U(1)$ 的情形。让我们在 Hamilton 框架中简要地讨论四维阿贝尔对偶。
取形如 $M_{3}\times\mathbb{R}$ 的四维流形，其中 $\mathbb{R}$ 是类时方向。
注意这是一个自旋流形，所以我们预期有完整的 $SL(2,\mathbb{Z})$ 对称性。
为简单起见，我们假定 $H_{1}(M_{3})$ 中没有挠。
每个类 $x\in H^{2}(M_{3})$ 决定了三维流形 $M_{3}$ 上的一个复线丛 $\mathcal{L}_{x}$（满足 $c_{1}(\mathcal{L}_{x})=x$）。
我们在三维流形 $M_{3}$ 上的理论，其 Hilbert 空间可以写成如下形式
\begin{equation}
  \tag{8.75}
  \label{eq:8-75}
  \mathcal{H}_{\tau}(M_{3})=\bigoplus_{x\in H^{2}(M_{3},\mathbb{Z})}\mathcal{H}_{x},
\end{equation}
其中 $\mathcal{H}_{x}$ 是把 $\mathcal{L}_{x}$ 上的联络量子化所得的 Hilbert 空间。
（在左端，我们显式标出了对耦合常数 $\tau$ 的依赖。）
为构造 $\mathcal{H}_{x}$，把任意联络写成 $A=A_{0}+\beta$ 的形式，其中 $A_{0}$ 是满足某个规范固定条件的调和联络（harmonic connection，即曲率为调和二形式的联络；规范固定的细节无关紧要），而 $\beta$ 是一个 1-形式。
令 $\mathcal{T}_{x}$ 为线丛 $\mathcal{L}_{x}$ 上调和联络的空间模去规范变换。于是量子化给出
\begin{equation}
  \tag{8.76}
  \label{eq:8-76}
  \mathcal{H}_{x}=\mathcal{H}_{\beta}\otimes L^{2}(\mathcal{T}_{x}).
\end{equation}

这里 $\mathcal{H}_{\beta}$ 是把 $\beta$ 的空间量子化所得的 Hilbert 空间，而 $L^{2}(\mathcal{T}_{x})$ 就是 $\mathcal{T}_{x}$ 上的 $L^{2}$ 函数空间。

注意，因子 $\mathcal{H}_{\beta}$ 与 $x$ 无关，因为 1-形式的空间在定义时并未涉及 $x$。
事实上，$\mathcal{H}_{\beta}$ 是 $M_{3}$ 上一个阿贝尔规范理论的 Hilbert 空间，只不过该理论的规范群不是 $U(1)$，而是它的万有覆盖 $\mathbb{R}$。
对偶把 $\mathcal{H}_{\beta}$ 映到自身，同时分别作用在 $\mathcal{H}'=\bigoplus_{x}L^{2}(\mathcal{T}_{x})$ 上。
对偶在 $\mathcal{H}_{\beta}$ 上的作用由 (8.74) 中的算符映射给出。

对偶在 $\mathcal{H}'$ 上的作用可以这样描述。
注意，所有的 $\mathcal{T}_{x}$ 都是环面 $H^{1}(M_{3},\mathbb{R}/\mathbb{Z})$ 作用下的主齐性空间（principal homogeneous space），而该环面参数化了 $M_{3}$ 上的平坦线丛；$H^{1}(M_{3},\mathbb{R}/\mathbb{Z})$ 中一个元素的作用，就是把给定的线丛（带联络）与相应的平坦线丛作张量积。
令 $y$ 表示阿贝尔群 $H^{1}(M_{3},\mathbb{R}/\mathbb{Z})$ 的一个特征标（character）。存在一个分解
% 原文在此处跨页：句子于 p.126 续完（"decomposition $L^2(\mathcal T_x)=\bigoplus_y\mathcal T_{x,y}$"）。

% [Tutor Note: 这里 $\mathcal{T}_{x}$ 的几何值得写清楚，否则 (8.76)、(8.77) 的分解会显得神秘。两个调和联络之差是闭 1-形式模恰当形式，即 $H^{1}(M_{3},\mathbb{R})$；固定 $c_{1}=x$ 只固定了「拓扑」部分，故 $\mathcal{T}_{x}$ 是一个维数为 $b_{1}$ 的环面，而且是 $H^{1}(M_{3},\mathbb{R}/\mathbb{Z})$ 的主齐性空间（torsor）而非线性空间——(8.75) 左端的 $\tau$ 依赖恰好体现在「哪一个 torsor」上：不同的平坦丛扭转不改变 $L^{2}(\mathcal{T}_x)$ 的 Hilbert 空间结构，只改变其上的平坦联络，这正是 (8.80) 中出现平坦线丛 $\mathscr L_{\theta}$ 的由来。另注意 $\mathcal{H}_{\beta}$ 对应「规范群为 $\mathbb{R}$」而非 $U(1)$：因为 $\beta$ 描述的是 $A$ 偏离调和代表的部分，其「大规范变换」不受整性约束，故其量子化给出的是 $\mathbb{R}$ 规范理论的 Hilbert 空间；这也解释了为什么对偶只作用在 $\mathcal{H}'$ 上（见 (8.74)：$F_\pm$ 的插入只涉及 $F_C$，不涉及 $\beta$）。]

% ===== 原文 p.126 =====
$L^{2}(\mathcal{T}_{x})=\bigoplus_{y}\mathcal{T}_{x,y}$，其中 $\mathcal{T}_{x,y}$ 是 $L^{2}(\mathcal{T}_{x})$ 中按特征标 $y$ 变换的子空间。
每个 $\mathcal{T}_{x,y}$ 都是一维的。因此
\begin{equation}
  \tag{8.77}
  \label{eq:8-77}
  \mathcal{H}'=\bigoplus_{x,y}\mathcal{T}_{x,y}.
\end{equation}

注意，由 Poincaré 对偶与 Pontryagin 对偶，$H^{1}(M_{3},\mathbb{R}/\mathbb{Z})$ 的特征标群是 $H^{2}(M_{3},\mathbb{Z})$。
于是 $x$ 与 $y$ 取值于同一个空间。因此可以交换它们，而这正是 $\tau\to-1/\tau$ 变换所做的（更确切地说，它按 $(x,y)\to(-y,x)$ 作用）。
因此，对偶把一个经典的概念——按 $x$ 所作的分解——与一个量子的概念——按 $y$ 所作的分解——交换了。
关于对偶如何作用的这一断言，将在下面引入算符 $Q_{e}$ 与 $Q_{m}$ 的地方得到验证。

% [Tutor Note: 「由 Poincaré 与 Pontryagin 对偶，$H^{1}(M_{3},\mathbb{R}/\mathbb{Z})$ 的特征标群是 $H^{2}(M_{3},\mathbb{Z})$」这一步，用的正是 p.125 开头「假定 $H_{1}(M_{3})$ 无挠」这一前提，值得补全：$H^{1}(M,\mathbb{R}/\mathbb{Z})=\operatorname{Hom}(H_{1}(M,\mathbb{Z}),\mathbb{R}/\mathbb{Z})$ 的 Pontryagin 对偶是 $H_{1}(M,\mathbb{Z})$ 的\emph{自由部分}，再由 Poincaré 对偶 $H_{1}(M,\mathbb{Z})\cong H^{2}(M,\mathbb{Z})$（配合万有系数定理：$H^{2}$ 的挠部分恰等于 $H_{1}$ 的挠部分）才等于 $H^{2}(M,\mathbb{Z})$。若存在挠，$x$ 与 $y$ 就不再取值于同一集合，$(x,y)\to(-y,x)$ 的写法须改为对 $\operatorname{Hom}$ 的表述，$\mathcal{H}'=\bigoplus_{x,y}\mathcal{T}_{x,y}$ 也不再是同一个环面上的 Fourier 分解。这也说明作者为何在 p.125 一开始就把「无挠」作为简化假设单独提出。]

在量子化时——先暂时略去 $\theta$——人们把四维曲率写成 $F_{A}'+e^{2}\pi_{A}dt$，其中 $F_{A}'$ 是 $M_{3}$ 上的一个二形式，而 $\pi_{A}$——$M_{3}$ 上的一个一形式——是联络 $A$ 的共轭动量。Hamilton 量成为
\begin{equation}
  \tag{8.78}
  \label{eq:8-78}
  H=\int d^{3}x\left(\frac{1}{2e^{2}}F_{A_{0}}^{2}+\frac{e^{2}}{2}\nabla_{A_{0}}^{2}\right)+H(\beta).
\end{equation}

这里 $H(\beta)$ 是 Hamilton 量中作用在 $\mathcal{H}_{\beta}$ 上的那一部分。其余各项作用在 $\mathcal{H}'$ 上。
第一项是调和联络 $A_{0}$ 的磁能；它来自拉氏量中关于 $F_{A}'$ 二次的部分，是 $\int_{M_{3}}x\wedge *x$ 的一个倍数。
第二项来自拉氏量中关于 $\pi_{A}$ 二次的部分，是电能，即 $\mathcal{T}_{x}$ 上的 Laplacian $\nabla^{2}$；它是 $\int_{M_{3}}y\wedge *y$ 的一个倍数。

把 $\theta$ 项包括进来会使量子化发生改变。
事实上，正则动量 $\pi_{A}$——出于后面将显出的理由，我们把它写成 $F_{A}^{\vee}=*\pi_{A}$——可以由原来的拉氏量 (8.58) 定出，其结果为
\begin{equation}
  \tag{8.79}
  \label{eq:8-79}
  F_{A}^{\vee}=2\pi i\frac{\delta\mathcal{L}}{\delta F_{A}}=\frac{2\pi i}{e^{2}}*F_{A}-\frac{\theta}{\pi}F_{A}.
\end{equation}

在 $\theta\ne0$ 时，正则量子化给出的不是 $\mathcal{H}'=\bigoplus_{x}L^{2}(\mathcal{T}_{x})$，而是
\begin{equation}
  \tag{8.80}
  \label{eq:8-80}
  \mathcal{H}'=\bigoplus_{x}\Gamma_{L^{2}}(\mathcal{T}_{x},\mathscr{L}_{\theta}),
\end{equation}
其中 $\mathscr{L}_{\theta}$ 是 $\mathcal{T}_{x}$ 上某个平坦线丛，而 $\Gamma_{L^{2}}$ 表示 $L^{2}$ 截面的空间。
我把辨认 $\mathscr{L}_{\theta}$ 留作读者的一个练习。\footnote{做这个练习的出发点应当是注意到：在一个把联络在时间方向上平凡化的规范中，拉氏量中的 $\theta$ 项是齐次的、并且对时间导数是线性的；任何这样的项都会在相空间上引出一个带联络的线丛。在目前的情形，$\theta$ 项是一个拓扑不变量，从而是经典位形空间上的一个局部常值函数；这就导致线丛 $\mathscr{L}_{\theta}$ 是平坦的。}
当然，$\theta=0$ 时的 $\mathscr{L}_{\theta}$ 就是具有平凡平坦联络的平坦丛；而 $\mathscr{L}_{\theta}$ 只依赖于 $\theta$ 模 $\pi$——我们在自旋流形上，所以 $\theta\to\theta+\pi$ 是一个对称性。

把 $F_{A}^{\vee}$ 的公式用 $\tau$ 改写，我们就能确定这个算符在 $\tau\to-1/\tau$ 下如何变换。
事实上，在算符映射 (8.74) 之下，人们得到
\begin{equation}
  \tag{8.81}
  \label{eq:8-81}
  F_{A}^{\vee}=-\bar\tau F_{+}+\tau F_{-}\mapsto F_{C+}+F_{C-}=F_{C}.
\end{equation}

我们还需要这个计算的「对偶」版本：
\begin{equation}
  \tag{8.82}
  \label{eq:8-82}
  F_{A}=F_{+}+F_{-}\mapsto\overline{\left(-1/\tau\right)}F_{C+}-\left(-1/\tau\right)F_{C-}=-F_{C}^{\vee}.
\end{equation}

% [Tutor Note: (8.79) 与 (8.81) 不可能同时成立，这是本页最需要提防的符号陷阱。把 $\tau=\frac{\theta}{\pi}+\frac{2\pi i}{e^{2}}$、$F_{\pm}=\frac12(F_{A}\pm*F_{A})$ 代入 (8.79) 印出的右端 $\frac{2\pi i}{e^{2}}*F_{A}-\frac{\theta}{\pi}F_{A}$，得到 $-\bar\tau F_{+}-\tau F_{-}$；而 (8.81) 印的是 $-\bar\tau F_{+}+\tau F_{-}$，$F_{-}$ 项差一个符号。三式 (8.81)、(8.82)、(8.83) 构成自洽的一组：把 (8.81) 的形式记作 $F^{\vee}=-\bar\tau F_{+}+\tau F_{-}$，则对 $\tau'=-1/\tau$ 有 $F_{C}^{\vee}=-\bar\tau'F_{C+}+\tau'F_{C-}=\frac{1}{\bar\tau}F_{C+}-\frac{1}{\tau}F_{C-}$，于是 (8.82) 的右端 $\overline{(-1/\tau)}F_{C+}-(-1/\tau)F_{C-}=-\frac{1}{\bar\tau}F_{C+}+\frac{1}{\tau}F_{C-}=-F_{C}^{\vee}$ 恰好成立，(8.83) 的 $Q_{e}\mapsto Q_{m}$、$Q_{m}\mapsto-Q_{e}$ 也随之成立。故 (8.79) 中 $F_{A}$ 与 $*F_{A}$ 两处应互换，即应为 $F_{A}^{\vee}=\frac{2\pi i}{e^{2}}F_{A}-\frac{\theta}{\pi}*F_{A}$；它与 (8.81) 只差这一点。但要注意：p.127 上「$\theta$ 增加 $\pi$ 时 $Q_{e}\mapsto Q_{e}+Q_{m}$」（即 Witten 效应）这句话，用 (8.79) 印出的形式才给出正比于 $Q_{m}$ 的平移，用互换后的形式给出的却是正比于电通量的平移。可见原文此处至少有一处符号笔误；读者务必先固定一套约定（例如以 (8.82)、(8.83) 为准），再自行复核，切勿两式混用。]

% ===== 原文 p.127 =====
在 Hamilton 形式中，对任何 2-循环 $\Sigma\subset M_{3}$，我们都可以在 Hilbert 空间 $\mathcal{H}_{\tau}(M_{3})$ 上定义与之相应的「电」算符与「磁」算符，即令
\begin{equation}
  \tag{8.83}
  \label{eq:8-83}
  \begin{aligned}
  Q_{e}(\Sigma)&=\int_{\Sigma}\frac{F_{A}^{\vee}}{2\pi}=\int_{\Sigma}\frac{F_{C}}{2\pi},\\
  Q_{m}(\Sigma)&=\int_{\Sigma}\frac{F_{A}}{2\pi}=-\int_{\Sigma}\frac{F_{C}^{\vee}}{2\pi}.
  \end{aligned}
\end{equation}
（这些算符只依赖于 $\Sigma$ 在 $H_{2}(M_{3},\mathbb{Z})$ 中的类。）
显然，在 $\tau\to-1/\tau$ 之下，也就是在 $A\to C$ 之下，有 $Q_{e}\to Q_{m}$、$Q_{m}\to-Q_{e}$。
既然 $x$ 与 $y$ 分别是 $Q_{m}$ 与 $Q_{e}$ 的本征值，这就意味着 $\tau\to-1/\tau$ 按 $(x,y)\to(-y,x)$ 作用。
此外，由 $F_{A}^{\vee}$ 的显式公式 (8.79) 可以看出：当 $\theta$ 增加 $\pi$ 时，算符 $Q_{m}$ 不变，而算符 $Q_{e}$ 变为 $Q_{e}+Q_{m}$。
由于我们在自旋流形上，$\theta\to\theta+\pi$ 是一个对称性；它对应于 $\tau\to\tau+1$，并且不改变 Hilbert 空间。

上一段所作的陈述可以合并为如下一句：$SL(2,\mathbb{Z})$ 通过它在序对 $(x,y)$ 上的自然作用而作用在 $\mathcal{H}'$ 上。

\subsection*{参考文献与交叉引用}
% 原文：References and cross references

本讲的主要参考文献为
\begin{itemize}
  \item T. Buscher, \emph{Path integral derivation of quantum duality in nonlinear sigma models}, Phys. Lett. B \textbf{201} (1988), 466--472.
  \item E. Witten, \emph{On $S$-duality in abelian gauge theory}, Selecta Math. (N.S.) \textbf{1} (1995), 383--410.
  \item A. M. Polyakov, \emph{Gauge Fields and Strings}, Harwood Academic Publishers, 1987（特别是 pp.~62--72）.
  \item S. Coleman, \emph{Aspects of Symmetry}, Cambridge University Press, 1985（论「Lumps」的那一章）.
\end{itemize}
